% Include with % Include with % Include with % Include with \input{bankruptcy_balance_sheet.tex}.
% Requires \usepackage{tikz}.
\begin{tikzpicture}[
    font=\sffamily,
    scale=0.8,
    every node/.style={transform shape}
]

% Muted fill colors (editable).
\definecolor{cashfill}{HTML}{E3CB78}
\definecolor{debtfill}{HTML}{C98078}
\definecolor{unsecuredfill}{HTML}{DCA69E}
\definecolor{equityfill}{HTML}{7FA5C2}

% =========================
% Editable amounts (thousands of dollars).
% 2025 BAPCPA Table 1D: national totals / 181,587 complete schedules.
% Mortgage debt proxies all secured claims; home equity proxies real property
% minus secured claims. Unsecured debt uses nonpriority unsecured claims.
% https://www.uscourts.gov/sites/default/files/document/bapcpa_1d_1231.2025.pdf
% Cash: 3 * (2021 ABI monthly income 3,513 - monthly expenses 2,810).
% https://abi-org.s3.amazonaws.com/Covid19/stats/Weekly_Update_Dec_6-12_2021.pdf
% This stylized illustration mixes years; exempt equity is assumed to be 75%.
% Displayed dollar amounts and their box sizes are rounded to the nearest $10,000.
% =========================
\def\cash{2.109}
\def\homeequity{20}
\def\mortgagedebt{160}
\def\unsecureddebt{70}
\def\legalfees{2}
% Illustrative height for legal fees and stigma, including nonmonetary costs.
\def\legalcostheight{0.75}
\pgfmathsetmacro{\exemptassets}{0.75*\homeequity}
\newcommand{\amountlabel}[1]{(\$\pgfmathprintnumber[fixed,precision=0,1000 sep={,}]{#1e3})}

% Geometry: common height-per-unit scale before and after bankruptcy.
\def\heightperunit{0.025}
\def\barwidth{2.7}
\def\columngap{0.35}
\def\panelgap{1.0}
\pgfmathsetmacro{\debt}{\mortgagedebt+\unsecureddebt}
\pgfmathsetmacro{\top}{\heightperunit*max(\cash+\homeequity,\debt,\exemptassets)}
\pgfmathsetmacro{\panelwidth}{2*\barwidth+\columngap}
\pgfmathsetmacro{\panelshift}{\panelwidth+\panelgap}
\pgfmathsetmacro{\debtx}{\barwidth+\columngap}
\pgfmathsetmacro{\lostequity}{\homeequity-\exemptassets}
\pgfmathsetmacro{\netpanelshift}{2*\panelshift}
\pgfmathsetmacro{\netbaseline}{\top/2}

% =========================
% Left: Before bankruptcy
% =========================
\node at ({\panelwidth/2},{\top+1.1}) {\Large\textbf{Before bankruptcy}};
\node at ({\barwidth/2},{\top+0.45}) {\large Assets};
\node at ({\debtx+\barwidth/2},{\top+0.45}) {\large Liabilities};

% Assets: home equity below cash.
\draw[thick,fill=equityfill] (0,0) rectangle (\barwidth,{\heightperunit*\homeequity});
\node[align=center] at ({\barwidth/2},-0.65) {Home equity\\\amountlabel{\homeequity}};
\draw[thin] ({\barwidth/2},-0.2) -- ({\barwidth/2},0);
% Assumed exemption threshold: 75% of initial home equity.
\draw[red,very thick] (-0.08,{\heightperunit*\exemptassets})
    -- ({\barwidth+0.08},{\heightperunit*\exemptassets});
\node[red,align=center,anchor=east,font=\small]
    at (-0.18,{\heightperunit*\exemptassets}) {Homestead\\exemption};
\draw[thick,fill=cashfill] (0,{\heightperunit*\homeequity})
    rectangle (\barwidth,{\heightperunit*(\homeequity+\cash)});
% Place the cash label outside its proportionally small box.
\node[align=center,anchor=south] at ({\barwidth/2},{\heightperunit*(\homeequity+\cash)+0.35}) {Cash (?)};
\draw[thin] ({\barwidth/2},{\heightperunit*(\homeequity+\cash)})
    -- ({\barwidth/2},{\heightperunit*(\homeequity+\cash)+0.3});

% Liabilities: mortgage debt below unsecured debt.
\draw[thick,fill=debtfill] (\debtx,0)
    rectangle ({\debtx+\barwidth},{\heightperunit*\mortgagedebt});
\node[align=center] at ({\debtx+\barwidth/2},{\heightperunit*\mortgagedebt/2}) {Mortgage debt\\\amountlabel{\mortgagedebt}};
\draw[thick,fill=unsecuredfill] (\debtx,{\heightperunit*\mortgagedebt})
    rectangle ({\debtx+\barwidth},{\heightperunit*\debt});
\node[align=center] at ({\debtx+\barwidth/2},{\heightperunit*(\mortgagedebt+\unsecureddebt/2)}) {Unsecured debt\\\amountlabel{\unsecureddebt}};

\draw[dotted,thick]
    ({\panelwidth+\panelgap/2},-0.2) --
    ({\panelwidth+\panelgap/2},{\top+0.8});

% =========================
% Right: After bankruptcy
% =========================
\begin{scope}[xshift={\panelshift cm}]
    \node at ({\panelwidth/2},{\top+1.1}) {\Large\textbf{After bankruptcy}};
    \node at ({\barwidth/2},{\top+0.45}) {\large Assets};
    \node at ({\debtx+\barwidth/2},{\top+0.45}) {\large Liabilities};

    \draw[thick,fill=equityfill] (0,0) rectangle (\barwidth,{\heightperunit*\exemptassets});
    \node[align=center] at ({\barwidth/2},-0.65) {Homestead exemption\\(?)};
    \draw[thin] ({\barwidth/2},-0.2) -- ({\barwidth/2},0);

    % Mortgage debt remains; unsecured debt is discharged.
    \draw[thick,fill=debtfill] (\debtx,0)
        rectangle ({\debtx+\barwidth},{\heightperunit*\mortgagedebt});
    \node[align=center] at ({\debtx+\barwidth/2},{\heightperunit*\mortgagedebt/2}) {Mortgage debt\\\amountlabel{\mortgagedebt}};

\end{scope}

% Common baseline across both panels.
\draw[thick] (0,0) -- ({2*\panelwidth+\panelgap},0);

% =========================
% Third panel: gains and losses relative to a central zero line.
% Positive values are debt relief; negative values are legal fees and lost home equity.
% Use the same height-per-dollar scale as the first two panels.
% =========================
\draw[dotted,thick]
    ({2*\panelwidth+1.5*\panelgap},-0.2) --
    ({2*\panelwidth+1.5*\panelgap},{\top+0.8});
\begin{scope}[xshift={\netpanelshift cm}]
    \node at ({\panelwidth/2},{\top+1.1}) {\Large\textbf{Net change}};
    \pgfmathsetmacro{\netbarleft}{(\panelwidth-\barwidth)/2}
    \pgfmathsetmacro{\netbarright}{(\panelwidth+\barwidth)/2}
    \draw[thick,fill=unsecuredfill]
        (\netbarleft,\netbaseline)
        rectangle (\netbarright,{\netbaseline+\heightperunit*\unsecureddebt});
    \node[align=center] at
        ({\panelwidth/2},{\netbaseline+\heightperunit*\unsecureddebt/2})
        {Unsecured debt\\discharged\\\amountlabel{\unsecureddebt}};
    \draw[thick,fill=cashfill]
        (\netbarleft,\netbaseline)
        rectangle (\netbarright,{\netbaseline-\legalcostheight});
    \node[align=center,font=\small] at
        ({\panelwidth/2},{\netbaseline-\legalcostheight/2})
        {Fees and stigma};
    \draw[thick,fill=equityfill]
        (\netbarleft,{\netbaseline-\legalcostheight})
        rectangle (\netbarright,{\netbaseline-\legalcostheight-\heightperunit*\lostequity});
    \draw[thick] (0,\netbaseline) -- (\panelwidth,\netbaseline);
    \node[anchor=east] at (-0.1,\netbaseline) {$0$};
    \draw[thin] ({\panelwidth/2},{\netbaseline-\legalcostheight-\heightperunit*\lostequity})
        -- ({\panelwidth/2},{\netbaseline-\legalcostheight-\heightperunit*\lostequity-0.35});
    \node[align=center,anchor=north] at
        ({\panelwidth/2},{\netbaseline-\legalcostheight-\heightperunit*\lostequity-0.4})
        {Seized home equity\\(?)};
\end{scope}

\end{tikzpicture}
.
% Requires \usepackage{tikz}.
\begin{tikzpicture}[
    font=\sffamily,
    scale=0.8,
    every node/.style={transform shape}
]

% Muted fill colors (editable).
\definecolor{cashfill}{HTML}{E3CB78}
\definecolor{debtfill}{HTML}{C98078}
\definecolor{unsecuredfill}{HTML}{DCA69E}
\definecolor{equityfill}{HTML}{7FA5C2}

% =========================
% Editable amounts (thousands of dollars).
% 2025 BAPCPA Table 1D: national totals / 181,587 complete schedules.
% Mortgage debt proxies all secured claims; home equity proxies real property
% minus secured claims. Unsecured debt uses nonpriority unsecured claims.
% https://www.uscourts.gov/sites/default/files/document/bapcpa_1d_1231.2025.pdf
% Cash: 3 * (2021 ABI monthly income 3,513 - monthly expenses 2,810).
% https://abi-org.s3.amazonaws.com/Covid19/stats/Weekly_Update_Dec_6-12_2021.pdf
% This stylized illustration mixes years; exempt equity is assumed to be 75%.
% Displayed dollar amounts and their box sizes are rounded to the nearest $10,000.
% =========================
\def\cash{2.109}
\def\homeequity{20}
\def\mortgagedebt{160}
\def\unsecureddebt{70}
\def\legalfees{2}
% Illustrative height for legal fees and stigma, including nonmonetary costs.
\def\legalcostheight{0.75}
\pgfmathsetmacro{\exemptassets}{0.75*\homeequity}
\newcommand{\amountlabel}[1]{(\$\pgfmathprintnumber[fixed,precision=0,1000 sep={,}]{#1e3})}

% Geometry: common height-per-unit scale before and after bankruptcy.
\def\heightperunit{0.025}
\def\barwidth{2.7}
\def\columngap{0.35}
\def\panelgap{1.0}
\pgfmathsetmacro{\debt}{\mortgagedebt+\unsecureddebt}
\pgfmathsetmacro{\top}{\heightperunit*max(\cash+\homeequity,\debt,\exemptassets)}
\pgfmathsetmacro{\panelwidth}{2*\barwidth+\columngap}
\pgfmathsetmacro{\panelshift}{\panelwidth+\panelgap}
\pgfmathsetmacro{\debtx}{\barwidth+\columngap}
\pgfmathsetmacro{\lostequity}{\homeequity-\exemptassets}
\pgfmathsetmacro{\netpanelshift}{2*\panelshift}
\pgfmathsetmacro{\netbaseline}{\top/2}

% =========================
% Left: Before bankruptcy
% =========================
\node at ({\panelwidth/2},{\top+1.1}) {\Large\textbf{Before bankruptcy}};
\node at ({\barwidth/2},{\top+0.45}) {\large Assets};
\node at ({\debtx+\barwidth/2},{\top+0.45}) {\large Liabilities};

% Assets: home equity below cash.
\draw[thick,fill=equityfill] (0,0) rectangle (\barwidth,{\heightperunit*\homeequity});
\node[align=center] at ({\barwidth/2},-0.65) {Home equity\\\amountlabel{\homeequity}};
\draw[thin] ({\barwidth/2},-0.2) -- ({\barwidth/2},0);
% Assumed exemption threshold: 75% of initial home equity.
\draw[red,very thick] (-0.08,{\heightperunit*\exemptassets})
    -- ({\barwidth+0.08},{\heightperunit*\exemptassets});
\node[red,align=center,anchor=east,font=\small]
    at (-0.18,{\heightperunit*\exemptassets}) {Homestead\\exemption};
\draw[thick,fill=cashfill] (0,{\heightperunit*\homeequity})
    rectangle (\barwidth,{\heightperunit*(\homeequity+\cash)});
% Place the cash label outside its proportionally small box.
\node[align=center,anchor=south] at ({\barwidth/2},{\heightperunit*(\homeequity+\cash)+0.35}) {Cash (?)};
\draw[thin] ({\barwidth/2},{\heightperunit*(\homeequity+\cash)})
    -- ({\barwidth/2},{\heightperunit*(\homeequity+\cash)+0.3});

% Liabilities: mortgage debt below unsecured debt.
\draw[thick,fill=debtfill] (\debtx,0)
    rectangle ({\debtx+\barwidth},{\heightperunit*\mortgagedebt});
\node[align=center] at ({\debtx+\barwidth/2},{\heightperunit*\mortgagedebt/2}) {Mortgage debt\\\amountlabel{\mortgagedebt}};
\draw[thick,fill=unsecuredfill] (\debtx,{\heightperunit*\mortgagedebt})
    rectangle ({\debtx+\barwidth},{\heightperunit*\debt});
\node[align=center] at ({\debtx+\barwidth/2},{\heightperunit*(\mortgagedebt+\unsecureddebt/2)}) {Unsecured debt\\\amountlabel{\unsecureddebt}};

\draw[dotted,thick]
    ({\panelwidth+\panelgap/2},-0.2) --
    ({\panelwidth+\panelgap/2},{\top+0.8});

% =========================
% Right: After bankruptcy
% =========================
\begin{scope}[xshift={\panelshift cm}]
    \node at ({\panelwidth/2},{\top+1.1}) {\Large\textbf{After bankruptcy}};
    \node at ({\barwidth/2},{\top+0.45}) {\large Assets};
    \node at ({\debtx+\barwidth/2},{\top+0.45}) {\large Liabilities};

    \draw[thick,fill=equityfill] (0,0) rectangle (\barwidth,{\heightperunit*\exemptassets});
    \node[align=center] at ({\barwidth/2},-0.65) {Homestead exemption\\(?)};
    \draw[thin] ({\barwidth/2},-0.2) -- ({\barwidth/2},0);

    % Mortgage debt remains; unsecured debt is discharged.
    \draw[thick,fill=debtfill] (\debtx,0)
        rectangle ({\debtx+\barwidth},{\heightperunit*\mortgagedebt});
    \node[align=center] at ({\debtx+\barwidth/2},{\heightperunit*\mortgagedebt/2}) {Mortgage debt\\\amountlabel{\mortgagedebt}};

\end{scope}

% Common baseline across both panels.
\draw[thick] (0,0) -- ({2*\panelwidth+\panelgap},0);

% =========================
% Third panel: gains and losses relative to a central zero line.
% Positive values are debt relief; negative values are legal fees and lost home equity.
% Use the same height-per-dollar scale as the first two panels.
% =========================
\draw[dotted,thick]
    ({2*\panelwidth+1.5*\panelgap},-0.2) --
    ({2*\panelwidth+1.5*\panelgap},{\top+0.8});
\begin{scope}[xshift={\netpanelshift cm}]
    \node at ({\panelwidth/2},{\top+1.1}) {\Large\textbf{Net change}};
    \pgfmathsetmacro{\netbarleft}{(\panelwidth-\barwidth)/2}
    \pgfmathsetmacro{\netbarright}{(\panelwidth+\barwidth)/2}
    \draw[thick,fill=unsecuredfill]
        (\netbarleft,\netbaseline)
        rectangle (\netbarright,{\netbaseline+\heightperunit*\unsecureddebt});
    \node[align=center] at
        ({\panelwidth/2},{\netbaseline+\heightperunit*\unsecureddebt/2})
        {Unsecured debt\\discharged\\\amountlabel{\unsecureddebt}};
    \draw[thick,fill=cashfill]
        (\netbarleft,\netbaseline)
        rectangle (\netbarright,{\netbaseline-\legalcostheight});
    \node[align=center,font=\small] at
        ({\panelwidth/2},{\netbaseline-\legalcostheight/2})
        {Fees and stigma};
    \draw[thick,fill=equityfill]
        (\netbarleft,{\netbaseline-\legalcostheight})
        rectangle (\netbarright,{\netbaseline-\legalcostheight-\heightperunit*\lostequity});
    \draw[thick] (0,\netbaseline) -- (\panelwidth,\netbaseline);
    \node[anchor=east] at (-0.1,\netbaseline) {$0$};
    \draw[thin] ({\panelwidth/2},{\netbaseline-\legalcostheight-\heightperunit*\lostequity})
        -- ({\panelwidth/2},{\netbaseline-\legalcostheight-\heightperunit*\lostequity-0.35});
    \node[align=center,anchor=north] at
        ({\panelwidth/2},{\netbaseline-\legalcostheight-\heightperunit*\lostequity-0.4})
        {Seized home equity\\(?)};
\end{scope}

\end{tikzpicture}
.
% Requires \usepackage{tikz}.
\begin{tikzpicture}[
    font=\sffamily,
    scale=0.8,
    every node/.style={transform shape}
]

% Muted fill colors (editable).
\definecolor{cashfill}{HTML}{E3CB78}
\definecolor{debtfill}{HTML}{C98078}
\definecolor{unsecuredfill}{HTML}{DCA69E}
\definecolor{equityfill}{HTML}{7FA5C2}

% =========================
% Editable amounts (thousands of dollars).
% 2025 BAPCPA Table 1D: national totals / 181,587 complete schedules.
% Mortgage debt proxies all secured claims; home equity proxies real property
% minus secured claims. Unsecured debt uses nonpriority unsecured claims.
% https://www.uscourts.gov/sites/default/files/document/bapcpa_1d_1231.2025.pdf
% Cash: 3 * (2021 ABI monthly income 3,513 - monthly expenses 2,810).
% https://abi-org.s3.amazonaws.com/Covid19/stats/Weekly_Update_Dec_6-12_2021.pdf
% This stylized illustration mixes years; exempt equity is assumed to be 75%.
% Displayed dollar amounts and their box sizes are rounded to the nearest $10,000.
% =========================
\def\cash{2.109}
\def\homeequity{20}
\def\mortgagedebt{160}
\def\unsecureddebt{70}
\def\legalfees{2}
% Illustrative height for legal fees and stigma, including nonmonetary costs.
\def\legalcostheight{0.75}
\pgfmathsetmacro{\exemptassets}{0.75*\homeequity}
\newcommand{\amountlabel}[1]{(\$\pgfmathprintnumber[fixed,precision=0,1000 sep={,}]{#1e3})}

% Geometry: common height-per-unit scale before and after bankruptcy.
\def\heightperunit{0.025}
\def\barwidth{2.7}
\def\columngap{0.35}
\def\panelgap{1.0}
\pgfmathsetmacro{\debt}{\mortgagedebt+\unsecureddebt}
\pgfmathsetmacro{\top}{\heightperunit*max(\cash+\homeequity,\debt,\exemptassets)}
\pgfmathsetmacro{\panelwidth}{2*\barwidth+\columngap}
\pgfmathsetmacro{\panelshift}{\panelwidth+\panelgap}
\pgfmathsetmacro{\debtx}{\barwidth+\columngap}
\pgfmathsetmacro{\lostequity}{\homeequity-\exemptassets}
\pgfmathsetmacro{\netpanelshift}{2*\panelshift}
\pgfmathsetmacro{\netbaseline}{\top/2}

% =========================
% Left: Before bankruptcy
% =========================
\node at ({\panelwidth/2},{\top+1.1}) {\Large\textbf{Before bankruptcy}};
\node at ({\barwidth/2},{\top+0.45}) {\large Assets};
\node at ({\debtx+\barwidth/2},{\top+0.45}) {\large Liabilities};

% Assets: home equity below cash.
\draw[thick,fill=equityfill] (0,0) rectangle (\barwidth,{\heightperunit*\homeequity});
\node[align=center] at ({\barwidth/2},-0.65) {Home equity\\\amountlabel{\homeequity}};
\draw[thin] ({\barwidth/2},-0.2) -- ({\barwidth/2},0);
% Assumed exemption threshold: 75% of initial home equity.
\draw[red,very thick] (-0.08,{\heightperunit*\exemptassets})
    -- ({\barwidth+0.08},{\heightperunit*\exemptassets});
\node[red,align=center,anchor=east,font=\small]
    at (-0.18,{\heightperunit*\exemptassets}) {Homestead\\exemption};
\draw[thick,fill=cashfill] (0,{\heightperunit*\homeequity})
    rectangle (\barwidth,{\heightperunit*(\homeequity+\cash)});
% Place the cash label outside its proportionally small box.
\node[align=center,anchor=south] at ({\barwidth/2},{\heightperunit*(\homeequity+\cash)+0.35}) {Cash (?)};
\draw[thin] ({\barwidth/2},{\heightperunit*(\homeequity+\cash)})
    -- ({\barwidth/2},{\heightperunit*(\homeequity+\cash)+0.3});

% Liabilities: mortgage debt below unsecured debt.
\draw[thick,fill=debtfill] (\debtx,0)
    rectangle ({\debtx+\barwidth},{\heightperunit*\mortgagedebt});
\node[align=center] at ({\debtx+\barwidth/2},{\heightperunit*\mortgagedebt/2}) {Mortgage debt\\\amountlabel{\mortgagedebt}};
\draw[thick,fill=unsecuredfill] (\debtx,{\heightperunit*\mortgagedebt})
    rectangle ({\debtx+\barwidth},{\heightperunit*\debt});
\node[align=center] at ({\debtx+\barwidth/2},{\heightperunit*(\mortgagedebt+\unsecureddebt/2)}) {Unsecured debt\\\amountlabel{\unsecureddebt}};

\draw[dotted,thick]
    ({\panelwidth+\panelgap/2},-0.2) --
    ({\panelwidth+\panelgap/2},{\top+0.8});

% =========================
% Right: After bankruptcy
% =========================
\begin{scope}[xshift={\panelshift cm}]
    \node at ({\panelwidth/2},{\top+1.1}) {\Large\textbf{After bankruptcy}};
    \node at ({\barwidth/2},{\top+0.45}) {\large Assets};
    \node at ({\debtx+\barwidth/2},{\top+0.45}) {\large Liabilities};

    \draw[thick,fill=equityfill] (0,0) rectangle (\barwidth,{\heightperunit*\exemptassets});
    \node[align=center] at ({\barwidth/2},-0.65) {Homestead exemption\\(?)};
    \draw[thin] ({\barwidth/2},-0.2) -- ({\barwidth/2},0);

    % Mortgage debt remains; unsecured debt is discharged.
    \draw[thick,fill=debtfill] (\debtx,0)
        rectangle ({\debtx+\barwidth},{\heightperunit*\mortgagedebt});
    \node[align=center] at ({\debtx+\barwidth/2},{\heightperunit*\mortgagedebt/2}) {Mortgage debt\\\amountlabel{\mortgagedebt}};

\end{scope}

% Common baseline across both panels.
\draw[thick] (0,0) -- ({2*\panelwidth+\panelgap},0);

% =========================
% Third panel: gains and losses relative to a central zero line.
% Positive values are debt relief; negative values are legal fees and lost home equity.
% Use the same height-per-dollar scale as the first two panels.
% =========================
\draw[dotted,thick]
    ({2*\panelwidth+1.5*\panelgap},-0.2) --
    ({2*\panelwidth+1.5*\panelgap},{\top+0.8});
\begin{scope}[xshift={\netpanelshift cm}]
    \node at ({\panelwidth/2},{\top+1.1}) {\Large\textbf{Net change}};
    \pgfmathsetmacro{\netbarleft}{(\panelwidth-\barwidth)/2}
    \pgfmathsetmacro{\netbarright}{(\panelwidth+\barwidth)/2}
    \draw[thick,fill=unsecuredfill]
        (\netbarleft,\netbaseline)
        rectangle (\netbarright,{\netbaseline+\heightperunit*\unsecureddebt});
    \node[align=center] at
        ({\panelwidth/2},{\netbaseline+\heightperunit*\unsecureddebt/2})
        {Unsecured debt\\discharged\\\amountlabel{\unsecureddebt}};
    \draw[thick,fill=cashfill]
        (\netbarleft,\netbaseline)
        rectangle (\netbarright,{\netbaseline-\legalcostheight});
    \node[align=center,font=\small] at
        ({\panelwidth/2},{\netbaseline-\legalcostheight/2})
        {Fees and stigma};
    \draw[thick,fill=equityfill]
        (\netbarleft,{\netbaseline-\legalcostheight})
        rectangle (\netbarright,{\netbaseline-\legalcostheight-\heightperunit*\lostequity});
    \draw[thick] (0,\netbaseline) -- (\panelwidth,\netbaseline);
    \node[anchor=east] at (-0.1,\netbaseline) {$0$};
    \draw[thin] ({\panelwidth/2},{\netbaseline-\legalcostheight-\heightperunit*\lostequity})
        -- ({\panelwidth/2},{\netbaseline-\legalcostheight-\heightperunit*\lostequity-0.35});
    \node[align=center,anchor=north] at
        ({\panelwidth/2},{\netbaseline-\legalcostheight-\heightperunit*\lostequity-0.4})
        {Seized home equity\\(?)};
\end{scope}

\end{tikzpicture}
.
% Requires \usepackage{tikz}.
\begin{tikzpicture}[
    font=\sffamily,
    scale=0.8,
    every node/.style={transform shape}
]

% Muted fill colors (editable).
\definecolor{cashfill}{HTML}{E3CB78}
\definecolor{debtfill}{HTML}{C98078}
\definecolor{unsecuredfill}{HTML}{DCA69E}
\definecolor{equityfill}{HTML}{7FA5C2}

% =========================
% Editable amounts (thousands of dollars).
% 2025 BAPCPA Table 1D: national totals / 181,587 complete schedules.
% Mortgage debt proxies all secured claims; home equity proxies real property
% minus secured claims. Unsecured debt uses nonpriority unsecured claims.
% https://www.uscourts.gov/sites/default/files/document/bapcpa_1d_1231.2025.pdf
% Cash: 3 * (2021 ABI monthly income 3,513 - monthly expenses 2,810).
% https://abi-org.s3.amazonaws.com/Covid19/stats/Weekly_Update_Dec_6-12_2021.pdf
% This stylized illustration mixes years; exempt equity is assumed to be 75%.
% Displayed dollar amounts and their box sizes are rounded to the nearest $10,000.
% =========================
\def\cash{2.109}
\def\homeequity{20}
\def\mortgagedebt{160}
\def\unsecureddebt{70}
\def\legalfees{2}
% Illustrative height for legal fees and stigma, including nonmonetary costs.
\def\legalcostheight{0.75}
\pgfmathsetmacro{\exemptassets}{0.75*\homeequity}
\newcommand{\amountlabel}[1]{(\$\pgfmathprintnumber[fixed,precision=0,1000 sep={,}]{#1e3})}

% Geometry: common height-per-unit scale before and after bankruptcy.
\def\heightperunit{0.025}
\def\barwidth{2.7}
\def\columngap{0.35}
\def\panelgap{1.0}
\pgfmathsetmacro{\debt}{\mortgagedebt+\unsecureddebt}
\pgfmathsetmacro{\top}{\heightperunit*max(\cash+\homeequity,\debt,\exemptassets)}
\pgfmathsetmacro{\panelwidth}{2*\barwidth+\columngap}
\pgfmathsetmacro{\panelshift}{\panelwidth+\panelgap}
\pgfmathsetmacro{\debtx}{\barwidth+\columngap}
\pgfmathsetmacro{\lostequity}{\homeequity-\exemptassets}
\pgfmathsetmacro{\netpanelshift}{2*\panelshift}
\pgfmathsetmacro{\netbaseline}{\top/2}

% =========================
% Left: Before bankruptcy
% =========================
\node at ({\panelwidth/2},{\top+1.1}) {\Large\textbf{Before bankruptcy}};
\node at ({\barwidth/2},{\top+0.45}) {\large Assets};
\node at ({\debtx+\barwidth/2},{\top+0.45}) {\large Liabilities};

% Assets: home equity below cash.
\draw[thick,fill=equityfill] (0,0) rectangle (\barwidth,{\heightperunit*\homeequity});
\node[align=center] at ({\barwidth/2},-0.65) {Home equity\\\amountlabel{\homeequity}};
\draw[thin] ({\barwidth/2},-0.2) -- ({\barwidth/2},0);
% Assumed exemption threshold: 75% of initial home equity.
\draw[red,very thick] (-0.08,{\heightperunit*\exemptassets})
    -- ({\barwidth+0.08},{\heightperunit*\exemptassets});
\node[red,align=center,anchor=east,font=\small]
    at (-0.18,{\heightperunit*\exemptassets}) {Homestead\\exemption};
\draw[thick,fill=cashfill] (0,{\heightperunit*\homeequity})
    rectangle (\barwidth,{\heightperunit*(\homeequity+\cash)});
% Place the cash label outside its proportionally small box.
\node[align=center,anchor=south] at ({\barwidth/2},{\heightperunit*(\homeequity+\cash)+0.35}) {Cash (?)};
\draw[thin] ({\barwidth/2},{\heightperunit*(\homeequity+\cash)})
    -- ({\barwidth/2},{\heightperunit*(\homeequity+\cash)+0.3});

% Liabilities: mortgage debt below unsecured debt.
\draw[thick,fill=debtfill] (\debtx,0)
    rectangle ({\debtx+\barwidth},{\heightperunit*\mortgagedebt});
\node[align=center] at ({\debtx+\barwidth/2},{\heightperunit*\mortgagedebt/2}) {Mortgage debt\\\amountlabel{\mortgagedebt}};
\draw[thick,fill=unsecuredfill] (\debtx,{\heightperunit*\mortgagedebt})
    rectangle ({\debtx+\barwidth},{\heightperunit*\debt});
\node[align=center] at ({\debtx+\barwidth/2},{\heightperunit*(\mortgagedebt+\unsecureddebt/2)}) {Unsecured debt\\\amountlabel{\unsecureddebt}};

\draw[dotted,thick]
    ({\panelwidth+\panelgap/2},-0.2) --
    ({\panelwidth+\panelgap/2},{\top+0.8});

% =========================
% Right: After bankruptcy
% =========================
\begin{scope}[xshift={\panelshift cm}]
    \node at ({\panelwidth/2},{\top+1.1}) {\Large\textbf{After bankruptcy}};
    \node at ({\barwidth/2},{\top+0.45}) {\large Assets};
    \node at ({\debtx+\barwidth/2},{\top+0.45}) {\large Liabilities};

    \draw[thick,fill=equityfill] (0,0) rectangle (\barwidth,{\heightperunit*\exemptassets});
    \node[align=center] at ({\barwidth/2},-0.65) {Homestead exemption\\(?)};
    \draw[thin] ({\barwidth/2},-0.2) -- ({\barwidth/2},0);

    % Mortgage debt remains; unsecured debt is discharged.
    \draw[thick,fill=debtfill] (\debtx,0)
        rectangle ({\debtx+\barwidth},{\heightperunit*\mortgagedebt});
    \node[align=center] at ({\debtx+\barwidth/2},{\heightperunit*\mortgagedebt/2}) {Mortgage debt\\\amountlabel{\mortgagedebt}};

\end{scope}

% Common baseline across both panels.
\draw[thick] (0,0) -- ({2*\panelwidth+\panelgap},0);

% =========================
% Third panel: gains and losses relative to a central zero line.
% Positive values are debt relief; negative values are legal fees and lost home equity.
% Use the same height-per-dollar scale as the first two panels.
% =========================
\draw[dotted,thick]
    ({2*\panelwidth+1.5*\panelgap},-0.2) --
    ({2*\panelwidth+1.5*\panelgap},{\top+0.8});
\begin{scope}[xshift={\netpanelshift cm}]
    \node at ({\panelwidth/2},{\top+1.1}) {\Large\textbf{Net change}};
    \pgfmathsetmacro{\netbarleft}{(\panelwidth-\barwidth)/2}
    \pgfmathsetmacro{\netbarright}{(\panelwidth+\barwidth)/2}
    \draw[thick,fill=unsecuredfill]
        (\netbarleft,\netbaseline)
        rectangle (\netbarright,{\netbaseline+\heightperunit*\unsecureddebt});
    \node[align=center] at
        ({\panelwidth/2},{\netbaseline+\heightperunit*\unsecureddebt/2})
        {Unsecured debt\\discharged\\\amountlabel{\unsecureddebt}};
    \draw[thick,fill=cashfill]
        (\netbarleft,\netbaseline)
        rectangle (\netbarright,{\netbaseline-\legalcostheight});
    \node[align=center,font=\small] at
        ({\panelwidth/2},{\netbaseline-\legalcostheight/2})
        {Fees and stigma};
    \draw[thick,fill=equityfill]
        (\netbarleft,{\netbaseline-\legalcostheight})
        rectangle (\netbarright,{\netbaseline-\legalcostheight-\heightperunit*\lostequity});
    \draw[thick] (0,\netbaseline) -- (\panelwidth,\netbaseline);
    \node[anchor=east] at (-0.1,\netbaseline) {$0$};
    \draw[thin] ({\panelwidth/2},{\netbaseline-\legalcostheight-\heightperunit*\lostequity})
        -- ({\panelwidth/2},{\netbaseline-\legalcostheight-\heightperunit*\lostequity-0.35});
    \node[align=center,anchor=north] at
        ({\panelwidth/2},{\netbaseline-\legalcostheight-\heightperunit*\lostequity-0.4})
        {Seized home equity\\(?)};
\end{scope}

\end{tikzpicture}
